\documentclass[11pt,a4paper]{article}
\usepackage[utf8]{inputenc}
\usepackage[T1]{fontenc}
\usepackage{amsfonts}
\usepackage[french]{babel}
\frenchbsetup{StandardLists=true}
\usepackage{enumitem}
\usepackage{amssymb}
\def\nbR{\ensuremath{\mathrm{I\! R}}}
\usepackage{amsmath}
\usepackage{awesomebox}
\DeclareMathOperator{\e}{e}

\usepackage[left=2cm,right=2cm,top=2cm,bottom=2cm]{geometry}
\usepackage{multirow}
\usepackage[dvipsnames]{xcolor}
\usepackage{parskip}
\usepackage{caption}

\usepackage[T1]{fontenc}
\usepackage{fouriernc}
\usepackage{graphicx}
\usepackage{setspace}
\singlespacing
\usepackage{tcolorbox}
\usepackage{framed}
\usepackage{multicol}

\usepackage{fancybox}
\usepackage{fancyhdr}
\setlength{\headheight}{14pt}
\pagestyle{fancy}
\renewcommand\headrulewidth{1pt}
\fancyhead[L]{Lycée Denis Diderot}
\fancyhead[R]{\textcolor{red}{\textbf{VERSION PROFESSEUR}} -- 2BTS}
\setcounter{page}{1}\fancyfoot[C]{page \thepage}
\fancyfoot[R]{\today}
\fancyfoot[L]{Alexis RUIZ}

\usepackage{tikz}
\usepackage{pgfplots}
\pgfplotsset{compat=1.18}
\usepackage{tkz-tab}
\usetikzlibrary{arrows}
\usepackage{pifont}

\usepackage{multirow,tabularx,booktabs}
\usepackage{colortbl}
\usepackage[np]{numprint}
\usepackage[tikz]{bclogo}

\def\N{{\mathbb N}}
\def\Z{{\mathbb Z}}
\def\R{{\mathbb R}}
\def\C{{\mathbb C}}
\def\e{{\text{e}}}
\def\dt{{\text{d}t}}

\newcommand{\ud}{\mathrm{d}}

% Boîte de correction bleue
\newcommand{\cor}[1]{%
  \begin{tcolorbox}[
    colback=blue!5!white,
    colframe=blue!70!black,
    boxrule=1pt,
    arc=3pt,
    left=4pt, right=4pt, top=3pt, bottom=3pt,
    before skip=2mm, after skip=3mm
  ]
  \color{blue!80!black}#1
  \end{tcolorbox}%
}

% Boîte de résultat (vert)
\newcommand{\res}[1]{%
  \begin{tcolorbox}[
    colback=green!5!white,
    colframe=green!50!black,
    boxrule=0.8pt,
    arc=3pt,
    left=4pt, right=4pt, top=2pt, bottom=2pt,
    before skip=1mm, after skip=2mm
  ]
  \color{green!40!black}#1
  \end{tcolorbox}%
}

\newcounter{numexos}
\setcounter{numexos}{0}
\newcommand{\bex}[2]{%
\def\bextitle{#1}%
\addtocounter{numexos}{1}%
\begin{tcolorbox}[colback=white,colframe=black!30,arc=3pt,boxrule=0.6pt,
  left=4mm,right=4mm,top=1mm,bottom=2mm]%
  \textbf{Exercice~\thenumexos\ifx\bextitle\empty\else~:~~#1\fi}\par
  #2%
\end{tcolorbox}}

\newcommand{\cadreblanc}[1]{%
\medskip\framebox[\linewidth][c]{%
\rule{0mm}{#1mm}}\par
\medskip}

\setlength{\columnseprule}{.25mm}
\newcolumntype{C}[1]{>{\centering\arraybackslash}p{#1cm}}
\renewcommand{\arraystretch}{1.4}

\newtcolorbox{mybox}{colback=white!10!white,colframe=black!5!black}

\newcommand{\bdef}[1]
{\begin{bclogo}[couleur = white,logo=\bcinfo , arrondi = 0.1]
{~Définition~:}
~~
{#1}
~~
\end{bclogo}}

\newcommand{\rem}[1]
{\begin{bclogo}[couleur = white,logo=\bcinfo , arrondi = 0.1]
{~Remarque~:}
~~
{#1}
~~
\end{bclogo}}

\newcommand{\bpre}[1]
{\begin{bclogo}[couleur = white,logo=, barre=none, arrondi = 0.1]
{~Prérequis~:}
~~
{#1}
~~
\end{bclogo}}

%%%%%%%%%%%%%%%%%%%%%%%%%%%%%%%%%%%%%%%%%%%%%%%%%%%%%%%%%%%%%
\begin{document}

\begin{mybox}
\begin{center}
\textbf{\LARGE CHAPITRE 10 -- Séries de Fourier}
\end{center}
\end{mybox}

\vfill 

\begin{center}
\begin{minipage}{14cm}
\begin{bclogo}[couleur = white, arrondi = 0.1]
{~Capacités attendues~:}
~~
\begin{itemize}
\item Calculer les coefficients de Fourier dans des cas simples.
\item Savoir identifier le bon développement en série de Fourier d'une fonction donnée.
\item Savoir calculer la valeur approchée d'une valeur efficace d'une fonction en utilisant son développement en série de Fourier.
\end{itemize}
\end{bclogo}
\end{minipage}
\end{center}

\vfill 


\bpre{%
\setstretch{0.85}
On considère la fonction $f$ définie sur $\mathbb{R}$ par :
\begin{center}
$\begin{cases}
f(x) &\textrm{est paire}\\
f(x)&\textrm {est de période }2\pi \\
f(x)=1  & \textrm {pour~~} x \in \left[0;\dfrac{\pi}{2} \right]\\
f(x)=0  & \textrm {pour~~} x \in \left]\dfrac{\pi}{2};\pi \right]\\
\end{cases}$
\end{center}
\begin{enumerate}
    \item Représenter dans le repère ci-dessous la fonction $f$ sur l'intervalle $[-\pi;2\pi]$.
\end{enumerate}
}

\vfill 


\begin{center}
\begin{tikzpicture}[x=2cm, y=4cm]
\def\xmin{-2.3} \def\xmax{4.3} \def\ymin{-0.15} \def\ymax{1.25}
\draw[xstep=0.25,ystep=0.1,very thin,orange!25] (\xmin,\ymin) grid (\xmax,\ymax);
\draw[xstep=1,ystep=0.5,thin,orange!55] (\xmin,\ymin) grid (\xmax,\ymax);
\draw[->,>=stealth',thick,blue!70!black] (\xmin,0)--(\xmax,0) node[right]{$x$};
\draw[->,>=stealth',thick,blue!70!black] (0,\ymin)--(0,\ymax) node[above]{$y$};
\foreach \n/\lab in {-2/{-\pi},-1/{-\frac{\pi}{2}},1/{\frac{\pi}{2}},2/{\pi},3/{\frac{3\pi}{2}},4/{2\pi}}
    \draw[blue!70!black] (\n,2pt)--(\n,-6pt) node[below]{\small$\lab$};
\foreach \y in {0.5,1}
    \draw[blue!70!black] (2pt,\y)--(-5pt,\y) node[left]{\small$\y$};
% Courbe f : paire, T=2π — f=1 sur [-π/2 ; π/2] et [3π/2 ; 2π], f=0 ailleurs
\draw[very thick, red!70!black]
    (-2,0)--(-1,0)--(-1,1)--(1,1)--(1,0)--(3,0)--(3,1)--(4,1);
\draw[black!60,thick] (\xmin,\ymin) rectangle (\xmax,\ymax);
\end{tikzpicture}
\end{center}

\vfill 


\bpre{%
\setstretch{0.85}
\begin{enumerate}
\setcounter{enumi}{1}
\item Tracer la courbe représentative de la fonction $g(x)=\dfrac{1}{2}+\dfrac{2}{\pi}\cos(x)-\dfrac{2}{3\pi}\cos(3x)$ \\[-2mm]
\item Tracer la courbe représentative de la fonction $h(x)=g(x)+\dfrac{2}{5\pi}\cos(5x)$ \\[-2mm]
\item  Tracer la courbe représentative de la fonction $k(x)=h(x)-\dfrac{2}{7\pi}\cos(7x)$\\[-2mm]
\item Conjecturer une fonction dont la représentation graphique s'approcherait encore plus de celle de la fonction $f$.\\[-2mm]
\dotfill
\end{enumerate}


\cor{ 
  $$l(x)=\tfrac{1}{2}+\tfrac{2}{\pi}\cos x-\tfrac{2}{3\pi}\cos 3x+\tfrac{2}{5\pi}\cos 5x-\tfrac{2}{7\pi}\cos            
7x+\tfrac{2}{9\pi}\cos 9x-\tfrac{2}{11\pi}\cos 11x$$

}
}
\vfill

%%%%%%%%%%%%%%%%%%%%%%%%%%%
\newpage
%%%%%%%%%%%%%%%%%%%%%%%%%%%

\begin{mybox}
    \textbf{1. Calculs des coefficients de Fourier}
\end{mybox}

\begin{mybox}
    \textbf{1.1. Les séries de Fourier}
\end{mybox}

\bdef{
\textcolor{blue}{Si $f$ est une fonction périodique de période $T$ et de pulsation $\omega=\dfrac{2\pi}{T}$, son \textbf{développement en série de Fourier} est :
\[S(x) = a_0 + \sum_{n=1}^{+\infty}\bigl(a_n\cos(n\omega x) + b_n\sin(n\omega x)\bigr)\]}
}

Les nombres $a_0$, $a_n$ et $b_n$ \textbf{sont les coefficients de Fourier} de la fonction $f$.

\vfill 


\textbf{Exemple :} Dans le prérequis, $g(x)=\dfrac{1}{2}+\dfrac{2}{\pi}\cos(x)-\dfrac{2}{3\pi}\cos(3x)$ est le début du développement de $f$.

\cor{
$a_0=\dfrac{1}{2}$,\quad $a_1=\dfrac{2}{\pi}$,\quad $a_2=0$,\quad $a_3=\dfrac{-2}{3\pi}$,\quad $a_4=0$,\quad $a_5=\dfrac{2}{5\pi}$\\[2mm]
.
}

\vfill 

\rem{ Si $T$ (en seconde , $s$) est la période d'une fonction $f$ alors $\omega$ est appelé sa pulsation (en radians par seconde , $rad.s^{-1}$) , on note alors la fréquence d'un signal $f=\dfrac{\omega}{2\pi}$. Son unité étant le Hertz ($Hz$).

}

\vfill 

\begin{mybox}
    \textbf{1.2. Calculs des coefficients}
\end{mybox}

Nous admettrons les résultats suivants : \fbox{$a_0=\dfrac{1}{T} \displaystyle \int_\alpha^{\alpha+T} f(t)\mathrm{d}t$}\\[3mm]
$a_0$ est en fait la \textbf{valeur moyenne} de la fonction $f$ sur une période.
\vspace{3mm}

\rem{Le réel $\alpha$ est choisi arbitrairement et le résultat de l'intégrale ne dépend pas de ce choix. On prendra le plus souvent $\alpha = 0$ ou $\alpha = \dfrac{-T}{2}.$

}

\vspace{0.3cm}

\begin{multicols}{2}
\begin{center}
\fbox{$a_n=\dfrac{2}{T} \displaystyle \int_\alpha^{\alpha+T} f(t)\cos(n\omega t)\mathrm{d}t$} \\
\fbox{$b_n=\dfrac{2}{T} \displaystyle \int_\alpha^{\alpha+T} f(t)\sin(n\omega t)\mathrm{d}t$} \\
\end{center}\end{multicols}

\vfill 

\newpage

%%%%%%%%%%%%%%%%%%%%%%%%%%%%%%%%%%%%%%%%%%%%%%%%%%%%%%%%%%%%%

\rem{Si la fonction $f$ est de période $2\pi$ alors on a $\omega = 1$ et :\\[3 mm]
$a_0=\dfrac{1}{2\pi} \displaystyle \int_\alpha^{\alpha+2\pi} f(t)\mathrm{d}t$,\hfill $a_n=\dfrac{1}{\pi} \displaystyle \int_\alpha^{\alpha+2\pi} f(t)\cos(n t)\mathrm{d}t$, \hfill $b_n=\dfrac{1}{\pi} \displaystyle \int_\alpha^{\alpha+2\pi} f(t)\sin(n t)\mathrm{d}t$.

}


\fcolorbox[gray]{0}{0.9}{\textbf{Activité 1. Calculer les coefficients de Fourier à la main.}}\\[2mm]

On étudie la fonction $f$, périodique de période $T=2\pi$, définie sur $[0\,;\,2\pi[$ par :
$f(x)=1$ si $x\in[0\,;\,\pi]$, $f(x)=0$ si $x\in\,]\pi\,;\,2\pi]$.

On note $\omega = \dfrac{2\pi}{T} = 1$.

\textbf{Calcul de $a_0$}

\cor{
$$a_0 = \frac{1}{T}\int_0^{T} f(x)\,\mathrm{d}x
     = \frac{1}{2\pi}\int_0^{\pi} 1\,\mathrm{d}x
     = \frac{1}{2\pi}\bigl[\,x\,\bigr]_0^{\pi}
     = \frac{\pi}{2\pi} = \boxed{\frac{1}{2}}$$
}

\textbf{Calcul de $a_n$}

\cor{
	\vspace{8mm}
$$a_n = \frac{2}{T}\int_0^{T} f(x)\cos(n\omega x)\,\mathrm{d}x
      = \frac{1}{\pi}\int_0^{\pi}\cos(nx)\,\mathrm{d}x
      = \frac{1}{\pi}\left[\frac{\sin(nx)}{n}\right]_0^{\pi}
      = \frac{\sin(n\pi) - \sin(0)}{n\pi} = \boxed{0} \quad \forall\, n\geqslant 1$$ 

En effet, $\sin(n\pi)=0$ pour tout entier $n$.

\vspace{8mm}
}

\textbf{Calcul de $b_n$}

\cor{
	\vspace{8mm}
	
$$b_n = \frac{2}{T}\int_0^{T} f(x)\sin(n\omega x)\,\mathrm{d}x
      = \frac{1}{\pi}\int_0^{\pi}\sin(nx)\,\mathrm{d}x
      = \frac{1}{\pi}\left[\frac{-\cos(nx)}{n}\right]_0^{\pi}
      = \frac{1-\cos(n\pi)}{n\pi}$$ \\[3mm]

\begin{itemize}
  \item Si $n$ est \textbf{pair} : $\cos(n\pi)=1$, donc $b_n = 0$.
  \item Si $n$ est \textbf{impair} : $\cos(n\pi)=-1$, donc $b_n = \dfrac{2}{n\pi}$.
\end{itemize}
\vspace{8mm}
}

\textbf{Expression de $S(x)$}

\cor{
	\vspace{8mm}
$$\boxed{S(x) = \frac{1}{2} + \frac{2}{\pi}\sum_{\substack{n=1\\n\text{ impair}}}^{+\infty}\frac{\sin(nx)}{n}
= \frac{1}{2} + \frac{2}{\pi}\left(\sin x + \frac{\sin 3x}{3} + \frac{\sin 5x}{5} + \cdots\right)}$$
\vspace{8mm}
}

\newpage
%%%%%%%%%%%%%%%%%%%%%%%%%%%%%%%%%%%%%%%%%%%%%%%%%%%%%%%%%%%%%
\fcolorbox[gray]{0}{0.9}{\textbf{Exercice 1}}\\[2mm]

On étudie $f$ périodique de période $T=2\pi$, avec $f=1$ sur $[0;\pi/2]$ et $f=0$ sur $]\pi/2;2\pi]$.

On note $\omega=1$.

\textbf{Calcul de $a_0$}

\cor{
$$a_0 = \frac{1}{2\pi}\int_0^{2\pi} f(x)\,\mathrm{d}x
= \frac{1}{2\pi}\int_0^{\pi/2}1\,\mathrm{d}x
= \frac{1}{2\pi}\cdot\frac{\pi}{2} = \boxed{\frac{1}{4}}$$
}

\textbf{Calcul de $a_n$}

\cor{
$$a_n = \frac{1}{\pi}\int_0^{\pi/2}\cos(nx)\,\mathrm{d}x
= \frac{1}{\pi}\left[\frac{\sin(nx)}{n}\right]_0^{\pi/2}
= \boxed{\frac{\sin\!\left(\dfrac{n\pi}{2}\right)}{n\pi}}$$

Remarque : $a_n=0$ si $n$ est pair (car $\sin(n\pi/2)=0$), $a_n=\pm\dfrac{1}{n\pi}$ si $n$ est impair.
}

\cor{
	\vspace{4mm}
Valeurs de $n=1$ à $n=5$ :
\begin{center}
\renewcommand{\arraystretch}{1.8}
\begin{tabular}{|>{\centering\arraybackslash}m{1.2cm}|*{5}{>{\centering\arraybackslash}m{2cm}|}}
\hline
$n$ & $1$ & $2$ & $3$ & $4$ & $5$ \\
\hline
$\sin\!\bigl(\tfrac{n\pi}{2}\bigr)$ & $1$ & $0$ & $-1$ & $0$ & $1$ \\
\hline
$a_n$ & $\dfrac{1}{\pi}$ & $0$ & $\dfrac{-1}{3\pi}$ & $0$ & $\dfrac{1}{5\pi}$ \\
\hline
\end{tabular}
\renewcommand{\arraystretch}{1.4}
\end{center}
\vspace{4mm}
}

\textbf{Calcul de $b_n$}

\cor{
$$b_n = \frac{1}{\pi}\int_0^{\pi/2}\sin(nx)\,\mathrm{d}x
= \frac{1}{\pi}\left[\frac{-\cos(nx)}{n}\right]_0^{\pi/2}
= \frac{1}{\pi}\cdot\frac{1-\cos\!\left(\dfrac{n\pi}{2}\right)}{n}
= \boxed{\frac{1-\cos\!\left(\dfrac{n\pi}{2}\right)}{n\pi}}$$
}

\cor{
	\vspace{8mm}
Valeurs de $n=1$ à $n=5$ :
\begin{center}
\renewcommand{\arraystretch}{1.8}
\begin{tabular}{|>{\centering\arraybackslash}m{3.2cm}|*{5}{>{\centering\arraybackslash}m{2cm}|}}
\hline
$n$ & $1$ & $2$ & $3$ & $4$ & $5$ \\
\hline
$1-\cos\!\bigl(\tfrac{n\pi}{2}\bigr)$ & $1$ & $2$ & $1$ & $0$ & $1$ \\
\hline
$b_n$ & $\dfrac{1}{\pi}$ & $\dfrac{1}{\pi}$ & $\dfrac{1}{3\pi}$ & $0$ & $\dfrac{1}{5\pi}$ \\
\hline
\end{tabular}
\renewcommand{\arraystretch}{1.4}
\end{center}
\vspace{8mm}
}

Écrire le développement en série de Fourier de la fonction f pour jusqu’à n = 4
\cor{
	\vspace{8mm}
$$S_4(x) = \frac{1}{4} + \frac{1}{\pi}\cos x + \frac{1}{\pi}\sin x + \frac{1}{\pi}\sin 2x - \frac{1}{3\pi}\cos 3x + \frac{1}{3\pi}\sin 3x$$
\vspace{8mm}
}

\newpage
%%%%%%%%%%%%%%%%%%%%%%%%%%%%%%%%%%%%%%%%%%%%%%%%%%%%%%%%%%%%%
\fcolorbox[gray]{0}{0.9}{\textbf{Exercice 2}}\\[2mm]

On étudie $f$ périodique de période $T=2\pi$, avec $f=1$ sur $[0;\pi]$ et $f=-1$ sur $]\pi;2\pi]$.

\textbf{Calcul de $a_0$}

\cor{
	\vspace{4mm}
$$a_0 = \frac{1}{2\pi}\int_0^{2\pi} f(x)\,\mathrm{d}x
= \frac{1}{2\pi}\left[\int_0^{\pi} 1\,\mathrm{d}x + \int_{\pi}^{2\pi}(-1)\,\mathrm{d}x\right]
= \frac{1}{2\pi}\bigl[\pi - \pi\bigr] = \boxed{0}$$
\vspace{4mm}
}

\textbf{Calcul de $a_n$}

\cor{
\begin{align*}
a_n &= \frac{1}{\pi}\int_0^{2\pi} f(x)\cos(nx)\,\mathrm{d}x \\
&= \frac{1}{\pi}\left[\int_0^{\pi}1\cdot\cos(nx)\,\mathrm{d}x + \int_{\pi}^{2\pi}(-1)\cos(nx)\,\mathrm{d}x\right] \\
&= \frac{1}{\pi}\left[\left[\frac{\sin(nx)}{n}\right]_0^{\pi} - \left[\frac{\sin(nx)}{n}\right]_{\pi}^{2\pi}\right] \\
&= \frac{1}{\pi}\left[\frac{\sin(n\pi)-\sin(0)}{n} - \frac{\sin(2n\pi)-\sin(n\pi)}{n}\right] \\
&= \frac{1}{\pi}\left[\frac{0-0}{n} - \frac{0-0}{n}\right] = \boxed{0}
\end{align*}
}

\textbf{Calcul de $b_n$}

\cor{
\begin{align*}
b_n &= \frac{1}{\pi}\int_0^{2\pi} f(x)\sin(nx)\,\mathrm{d}x \\
&= \frac{1}{\pi}\left[\int_0^{\pi}1\cdot\sin(nx)\,\mathrm{d}x + \int_{\pi}^{2\pi}(-1)\sin(nx)\,\mathrm{d}x\right] \\
&= \frac{1}{\pi}\left[\left[\frac{-\cos(nx)}{n}\right]_0^{\pi} - \left[\frac{-\cos(nx)}{n}\right]_{\pi}^{2\pi}\right] \\
&= \frac{1}{\pi}\left[\frac{1-\cos(n\pi)}{n} - \frac{\cos(n\pi)-1}{n}\right]
= \frac{2\bigl(1-\cos(n\pi)\bigr)}{n\pi}
\end{align*}

\begin{itemize}
  \item Si $n$ est \textbf{pair} : $\cos(n\pi)=1$, donc $b_n=0$.
  \item Si $n$ est \textbf{impair} : $\cos(n\pi)=-1$, donc $b_n=\dfrac{4}{n\pi}$.
\end{itemize}
}

\vspace{3mm}

\textbf{Expression de $S(x)$}

\cor{
$$\boxed{S(x) = \frac{4}{\pi}\sum_{\substack{n=1\\n\text{ impair}}}^{+\infty}\frac{\sin(nx)}{n}
= \frac{4}{\pi}\left(\sin x + \frac{\sin 3x}{3} + \frac{\sin 5x}{5}+\cdots\right)}$$
}

\vspace{3mm}

\textbf{Développement au rang 4}

\cor{
$$S_4(x) = \frac{4}{\pi}\sin x + \frac{4}{3\pi}\sin 3x$$
}

\vspace{5mm}

\newpage 


%%%%%%%%%%%%%%%%%%%%%%%%%%%%%%%%%%%%%%%%%%%%%%%%%%%%%%%%%%%%%
\fcolorbox[gray]{0}{0.9}{\textbf{Exercice 3}}\\[2mm]

On étudie $f$ périodique de période $T=2$, avec $f=1$ sur $[0;1[$ et $f=0$ sur $[1;2[$.

On note $\omega = \dfrac{2\pi}{T} = \pi$.

\textbf{Calcul de $a_0$}

\cor{
	\vspace{3mm}
$$a_0 = \frac{1}{T}\int_0^{T} f(x)\,\mathrm{d}x
= \frac{1}{2}\left[\int_0^{1}1\,\mathrm{d}x + \int_1^{2}0\,\mathrm{d}x\right]
= \frac{1}{2}\bigl[1 + 0\bigr] = \boxed{\frac{1}{2}}$$
\vspace{3mm}
}

\vfill 


\textbf{Calcul de $a_n$}

\cor{
	\vspace{3mm}
\begin{align*}
a_n &= \frac{2}{T}\int_0^{T} f(x)\cos(n\omega x)\,\mathrm{d}x \\
&= \int_0^{1}1\cdot\cos(n\pi x)\,\mathrm{d}x + \int_1^{2}0\cdot\cos(n\pi x)\,\mathrm{d}x \\
&= \left[\frac{\sin(n\pi x)}{n\pi}\right]_0^{1} \\
&= \frac{\sin(n\pi) - \sin(0)}{n\pi} = \frac{0}{n\pi} = \boxed{0}
\end{align*}
\vspace{3mm}

}

\vfill 

\textbf{Calcul de $b_n$}

\cor{
	\vspace{3mm}
	
\begin{align*}
b_n &= \frac{2}{T}\int_0^{T} f(x)\sin(n\omega x)\,\mathrm{d}x \\
&= \int_0^{1}1\cdot\sin(n\pi x)\,\mathrm{d}x + \int_1^{2}0\cdot\sin(n\pi x)\,\mathrm{d}x \\
&= \left[\frac{-\cos(n\pi x)}{n\pi}\right]_0^{1} \\
&= \frac{-\cos(n\pi)+\cos(0)}{n\pi} = \frac{1-\cos(n\pi)}{n\pi}
\end{align*}

\begin{itemize}
  \item Si $n$ est \textbf{pair} : $\cos(n\pi)=1$, donc $b_n=0$.
  \item Si $n$ est \textbf{impair} : $\cos(n\pi)=-1$, donc $b_n=\dfrac{2}{n\pi}$.
\end{itemize}
\vspace{3mm}

}

\vfill 


\cor{\vspace{3mm}
	
$$\boxed{S(x) = \frac{1}{2} + \frac{2}{\pi}\left(\sin(\pi x)+\frac{\sin(3\pi x)}{3}+\frac{\sin(5\pi x)}{5}+\cdots\right)}$$
\vspace{3mm}
}

\newpage
%%%%%%%%%%%%%%%%%%%%%%%%%%%%%%%%%%%%%%%%%%%%%%%%%%%%%%%%%%%%%
\begin{mybox}
    3. Calcul des coefficients de Fourier dans les cas particuliers
\end{mybox}

\begin{center}
\begin{minipage}{14cm}
\begin{bclogo}[couleur = white,logo=\bcetoile, arrondi = 0.1]{Bon à savoir}
\cor{
	\vspace{3mm}
\begin{itemize}
  \item Si $f$ est \textbf{impaire} : $\displaystyle\int_{-\alpha}^{\alpha} f(x)\,\mathrm{d}x = 0$
  \item Si $f$ est \textbf{paire} : $\displaystyle\int_{-\alpha}^{\alpha} f(x)\,\mathrm{d}x = 2\int_0^{\alpha} f(x)\,\mathrm{d}x$
\end{itemize}
\vspace{3mm}
}
\end{bclogo}
\end{minipage}
\end{center}

\vfill

\begin{mybox}
    3.1. Si la fonction $f$ est impaire
\end{mybox}

\vspace{3mm}

\cor{
\vspace{3mm}
Comme $f$ est impaire, tous les coefficients cosinus sont nuls :
$$a_0 = 0 \qquad \text{et} \qquad a_n = 0 \quad \forall\,n\geqslant 1$$
La formule des $b_n$ se simplifie (on intègre sur une demi-période) :
$$\boxed{b_n = \frac{4}{T}\int_0^{T/2} f(x)\sin(n\omega x)\,\mathrm{d}x}$$
\vspace{3mm}
}

\vfill

\begin{mybox}
    3.2. Si la fonction $f$ est paire
\end{mybox}

\vspace{3mm}

\cor{
	\vspace{3mm}
	
Comme $f$ est paire, tous les coefficients sinus sont nuls :
$$b_n = 0 \quad \forall\,n\geqslant 1$$
Les formules des $a_0$ et $a_n$ se simplifient (on intègre sur une demi-période) :
$$\boxed{a_0 = \frac{2}{T}\int_0^{T/2} f(x)\,\mathrm{d}x} \qquad \boxed{a_n = \frac{4}{T}\int_0^{T/2} f(x)\cos(n\omega x)\,\mathrm{d}x}$$
\vspace{3mm}
}

\vfill

\begin{center}
\begin{minipage}{14cm}
\begin{bclogo}[couleur = white,logo=\bcetoile, arrondi = 0.1]{Bon à savoir}
\cor{
Dans les deux cas, on dit qu'on \textbf{travaille sur une demi-période}.
}
\end{bclogo}
\end{minipage}
\end{center}

\vfill 


\newpage
%%%%%%%%%%%%%%%%%%%%%%%%%%%%%%%%%%%%%%%%%%%%%%%%%%%%%%%%%%%%%
\fcolorbox[gray]{0}{0.9}{\textbf{Activité A4 -- Cas particuliers (fonction paire)}}\\[2mm]

On étudie $f$ \textbf{paire}, périodique de période $T=\pi$, avec $f(x)=x$ sur $[0;\pi/2]$.

On note $\omega = \dfrac{2\pi}{T} = 2$.

\begin{center}
\begin{tikzpicture}[x=1.4cm, y=2.2cm]
\def\xmin{-4.4} \def\xmax{4.4} \def\ymin{-1.4} \def\ymax{1.4}
\draw[xstep=0.25,ystep=0.1,very thin,orange!25] (\xmin,\ymin) grid (\xmax,\ymax);
\draw[xstep=1,ystep=0.5,thin,orange!55] (\xmin,\ymin) grid (\xmax,\ymax);
\draw[->,>=stealth',thick,blue!70!black] (\xmin,0)--(\xmax,0) node[right]{$x$};
\draw[->,>=stealth',thick,blue!70!black] (0,\ymin)--(0,\ymax) node[above]{$y$};
\foreach \n/\lab in {-4/{-2\pi},-3/{-\frac{3\pi}{2}},-2/{-\pi},-1/{-\frac{\pi}{2}},1/{\frac{\pi}{2}},2/{\pi},3/{\frac{3\pi}{2}},4/{2\pi}}
    \draw[blue!70!black] (\n,2pt)--(\n,-6pt) node[below]{\small$\lab$};
\foreach \y in {-1,-0.5,0.5,1}
    \draw[blue!70!black] (2pt,\y)--(-5pt,\y) node[left]{\small$\y$};
\draw[very thick, red!70!black] (-4.4,0.4)--(-4,0)--(-3,1)--(-2,0)--(-1,1)--(0,0)--(1,1)--(2,0)--(3,1)--(4,0)--(4.4,0.4);
\draw[black!60,thick] (\xmin,\ymin) rectangle (\xmax,\ymax);
\end{tikzpicture}
\end{center}

\cor{
Comme $f$ est \textbf{paire} : $b_n=0$ pour tout $n\geqslant 1$.
}

\textbf{Calcul de $a_0$}

\cor{
Pour une fonction paire de période $T=\pi$ :
$$a_0 = \frac{2}{T}\int_0^{T/2} f(x)\,\mathrm{d}x
= \frac{2}{\pi}\int_0^{\pi/2} x\,\mathrm{d}x
= \frac{2}{\pi}\left[\frac{x^2}{2}\right]_0^{\pi/2}
= \frac{2}{\pi}\cdot\frac{\pi^2}{8} = \boxed{\frac{\pi}{4}}$$
}

\textbf{Calcul de $a_n$}

\cor{
$$a_n = \frac{4}{T}\int_0^{T/2} f(x)\cos(n\omega x)\,\mathrm{d}x
= \frac{4}{\pi}\int_0^{\pi/2} x\cos(2nx)\,\mathrm{d}x$$

On intègre par parties avec $u=x$ et $v'=\cos(2nx)$ :
$$\int_0^{\pi/2} x\cos(2nx)\,\mathrm{d}x
= \left[\frac{x\sin(2nx)}{2n}\right]_0^{\pi/2} - \int_0^{\pi/2}\frac{\sin(2nx)}{2n}\,\mathrm{d}x$$
$$= \frac{\pi}{2}\cdot\frac{\sin(n\pi)}{2n} - \frac{1}{2n}\left[\frac{-\cos(2nx)}{2n}\right]_0^{\pi/2}
= 0 + \frac{1}{4n^2}\bigl(\cos(n\pi)-1\bigr)
= \frac{(-1)^n-1}{4n^2}$$

Donc :
$$a_n = \frac{4}{\pi}\cdot\frac{(-1)^n-1}{4n^2} = \frac{(-1)^n-1}{\pi n^2}$$

\begin{itemize}
  \item Si $n$ est \textbf{pair} : $(-1)^n=1$, donc $a_n=0$.
  \item Si $n$ est \textbf{impair} : $(-1)^n=-1$, donc $a_n=\dfrac{-2}{\pi n^2}$.
\end{itemize}
}

\textbf{Expression de $S(x)$}

\cor{
$$\boxed{S(x) = \frac{\pi}{4} - \frac{2}{\pi}\sum_{\substack{n=1\\n\text{ impair}}}^{+\infty}\frac{\cos(2nx)}{n^2}
= \frac{\pi}{4} - \frac{2}{\pi}\left(\cos(2x)+\frac{\cos(6x)}{9}+\frac{\cos(10x)}{25}+\cdots\right)}$$
}

\newpage
%%%%%%%%%%%%%%%%%%%%%%%%%%%%%%%%%%%%%%%%%%%%%%%%%%%%%%%%%%%%%
\fcolorbox[gray]{0}{0.9}{\textbf{QCM -- Faire le point}}\\[2mm]

\cor{
\textbf{Question 1.} Si pour tout $n>0$, $a_n=0$, la fonction $f$ est $\ldots$

\textbf{Réponse : b) est impaire.}

En effet, pour une fonction impaire, tous les coefficients cosinus sont nuls ($a_n=0$ pour tout $n\geqslant 0$). La nullité des $a_n$ pour $n>0$ est caractéristique de l'absence de termes en cosinus dans le développement, ce qui est le signe d'une fonction impaire (ou nulle).
}

\vfill 


\cor{
\textbf{Question 2.} $f$ paire, $T=2\pi$, $f=\pi$ sur $[0;\pi/2[$ et $f=0$ sur $[\pi/2;\pi[$.

\textbf{Réponse : b) $a_n=\dfrac{2}{\pi}\displaystyle\int_0^{\pi/2}\pi\cos(nt)\,\mathrm{d}t$}

Car pour une fonction paire de période $T=2\pi$ ($\omega=1$) :
$$a_n = \frac{4}{T}\int_0^{T/2}f(x)\cos(nx)\,\mathrm{d}x = \frac{2}{\pi}\int_0^{\pi}f(x)\cos(nx)\,\mathrm{d}x$$
Et comme $f=0$ sur $[\pi/2;\pi]$, on n'intègre que sur $[0;\pi/2]$, d'où la réponse b).

L'option a) utilise $\frac{1}{\pi}$ (formule générale, non simplifiée pour les paires) et intègre sur $[0;\pi]$, ce qui doublerait le résultat.
L'option c) intègre sur $[0;\pi]$ avec le facteur $\frac{2}{\pi}$, ce qui serait correct si $f\neq 0$ sur tout $[0;\pi]$.
}

\vfill 


\cor{
\textbf{Question 3.} $f$ périodique de période $\pi$ (non précisée paire ou impaire), $f=\pi$ sur $[0;\pi/2[$ et $f=0$ sur $[\pi/2;\pi[$.

\textbf{Calcul effectif :} Avec $T=\pi$, $\omega=2$ :
\begin{align*}
a_0 &= \frac{1}{\pi}\int_0^{\pi}f(x)\,\mathrm{d}x = \frac{1}{\pi}\int_0^{\pi/2}\pi\,\mathrm{d}x = \frac{\pi}{2} \neq \pi \quad \Rightarrow\text{ a) faux}\\
a_n &= \frac{2}{\pi}\int_0^{\pi/2}\pi\cos(2nx)\,\mathrm{d}x = 2\left[\frac{\sin(2nx)}{2n}\right]_0^{\pi/2} = \frac{\sin(n\pi)}{n} = 0 \quad \forall n\geqslant 1\\
b_n &= \frac{2}{\pi}\int_0^{\pi/2}\pi\sin(2nx)\,\mathrm{d}x = \frac{1-\cos(n\pi)}{n} \neq 0 \text{ pour } n \text{ impair} \quad\Rightarrow\text{ b) faux}
\end{align*}
Aucune des trois propositions n'est parfaitement rédigée ; la réponse la plus proche est \textbf{c) (aucune règle simplificatrice ne s'applique directement)}, car $f$ n'est ni paire ni impaire et les simplifications de a) et b) sont fausses.
}

\vfill 


\cor{
\textbf{Question 4.} $f$ de période $2\pi$, $f(x)=x^2$ sur $[0;2\pi[$.

\textbf{Réponse : c) $f$ est ni paire ni impaire.}

En effet, $f(-1)=f(-1+2\pi)=f(2\pi-1)=(2\pi-1)^2\neq 1=f(1)$. Donc $f(-x)\neq f(x)$ en général, $f$ n'est pas paire. De même $f(-1)\neq -f(1)$, donc $f$ n'est pas impaire.
}

\vfill 


\cor{
\textbf{Question 5.} $T=4$, donc $\omega=\dfrac{2\pi}{T}=\dfrac{2\pi}{4}=\dfrac{\pi}{2}$.

\textbf{Réponse : a) $\omega=\dfrac{\pi}{2}$}
}

\vfill 


\newpage
%%%%%%%%%%%%%%%%%%%%%%%%%%%%%%%%%%%%%%%%%%%%%%%%%%%%%%%%%%%%%
\fcolorbox[gray]{0}{0.9}{\textbf{Exercice 4}}\\[2mm]

On étudie $f$ \textbf{impaire}, périodique de période $T=\pi$, avec $f(x)=\dfrac{1}{2}$ sur $\left[0;\dfrac{\pi}{4}\right[$
et $f(x)=0$ sur $\left]\dfrac{\pi}{4};\dfrac{\pi}{2}\right]$.

On note $\omega = \dfrac{2\pi}{T} = 2$.

\begin{center}
\begin{tikzpicture}[x=1.2cm, y=2.2cm]
\def\xmin{-4.4} \def\xmax{6.4} \def\ymin{-1.4} \def\ymax{1.4}
\draw[xstep=0.25,ystep=0.1,very thin,orange!25] (\xmin,\ymin) grid (\xmax,\ymax);
\draw[xstep=1,ystep=0.5,thin,orange!55] (\xmin,\ymin) grid (\xmax,\ymax);
\draw[->,>=stealth',thick,blue!70!black] (\xmin,0)--(\xmax,0) node[right]{$x$};
\draw[->,>=stealth',thick,blue!70!black] (0,\ymin)--(0,\ymax) node[above]{$y$};
\foreach \n/\lab in {-4/{-2\pi},-3/{-\frac{3\pi}{2}},-2/{-\pi},-1/{-\frac{\pi}{2}},1/{\frac{\pi}{2}},2/{\pi},3/{\frac{3\pi}{2}},4/{2\pi},5/{\frac{5\pi}{2}},6/{3\pi}}
    \draw[blue!70!black] (\n,2pt)--(\n,-6pt) node[below]{\small$\lab$};
\foreach \y in {-1,-0.5,0.5,1}
    \draw[blue!70!black] (2pt,\y)--(-5pt,\y) node[left]{\small$\y$};
\begin{scope}
\clip (\xmin,\ymin) rectangle (\xmax,\ymax);
% Saut -0.5→+0.5 en x = -4,-2,0,2,4,6  (période T=2 unités TikZ)
\foreach \s in {-4,-2,0,2,4,6} {
    \draw[very thick, red!70!black] ({\s-0.5},-0.5)--(\s,-0.5);   % f=-0.5
    \draw[very thick, red!70!black] (\s,0.5)--({\s+0.5},0.5);     % f=+0.5
    \draw[dashed, red!70!black] (\s,-0.5)--(\s,0.5);              % saut vertical
    \draw[dashed, red!70!black] ({\s-0.5},-0.5)--({\s-0.5},0);   % chute 0→-0.5
    \draw[dashed, red!70!black] ({\s+0.5},0)--({\s+0.5},0.5);    % chute 0.5→0
}
\foreach \s in {-4,-2,0,2,4} {
    \draw[very thick, red!70!black] ({\s+0.5},0)--({\s+1.5},0);   % f=0
}
\end{scope}
\draw[black!60,thick] (\xmin,\ymin) rectangle (\xmax,\ymax);
\end{tikzpicture}
\end{center}

\vfill 


\cor{
Comme $f$ est \textbf{impaire} : $a_0 = 0$ et $a_n = 0$ pour tout $n\geqslant 1$.
}

\vfill 

\textbf{Calcul de $b_n$}

\cor{
$$b_n = \frac{4}{T}\int_0^{T/2} f(x)\sin(n\omega x)\,\mathrm{d}x
= \frac{4}{\pi}\int_0^{\pi/2} f(x)\sin(2nx)\,\mathrm{d}x
= \frac{4}{\pi}\int_0^{\pi/4} \frac{1}{2}\sin(2nx)\,\mathrm{d}x$$

$$= \frac{2}{\pi}\left[\frac{-\cos(2nx)}{2n}\right]_0^{\pi/4}
= \frac{2}{\pi}\cdot\frac{1}{2n}\Bigl(1-\cos\!\Bigl(\frac{n\pi}{2}\Bigr)\Bigr)
= \boxed{\frac{1}{n\pi}\left(1-\cos\!\left(\frac{n\pi}{2}\right)\right)}$$

Calcul des premiers termes :
\begin{center}
\begin{tabular}{|c|c|c|c|c|c|}
\hline
$n$ & 1 & 2 & 3 & 4 & 5 \\
\hline
$\cos(n\pi/2)$ & $0$ & $-1$ & $0$ & $1$ & $0$ \\
\hline
$b_n$ & $\dfrac{1}{\pi}$ & $\dfrac{1}{\pi}$ & $\dfrac{1}{3\pi}$ & $0$ & $\dfrac{1}{5\pi}$ \\
\hline
\end{tabular}
\end{center}
}

\vfill 

\textbf{Expression de $S(x)$}

\cor{
$$S(x) = \frac{1}{\pi}\sin(2x)+\frac{1}{\pi}\sin(4x)+\frac{1}{3\pi}\sin(6x)+0\cdot\sin(8x)+\frac{1}{5\pi}\sin(10x)+\cdots$$

$$= \frac{1}{\pi}\sum_{n=1}^{+\infty}\frac{1-\cos\!\left(\dfrac{n\pi}{2}\right)}{n}\sin(2nx)$$
}

\vfill 


\newpage
%%%%%%%%%%%%%%%%%%%%%%%%%%%%%%%%%%%%%%%%%%%%%%%%%%%%%%%%%%%%%
\fcolorbox[gray]{0}{0.9}{\textbf{4. La formule de Parseval et ses applications}}\\[3mm]

\fcolorbox[gray]{0}{0.9}{\textbf{4.1. La valeur efficace d'un signal.}}\\[3mm]

\begin{center}
\fbox{Si $f$ est une fonction périodique de période $T$, sa valeur efficace est :
$f_{\text{eff}}=\sqrt{\dfrac{1}{T}\displaystyle\int_\alpha^{\alpha+T}\bigl(f(x)\bigr)^2\,\mathrm{d}x}$}
\end{center}

\vspace{5mm}

Cette formule est souvent utilisée en sciences physiques pour calculer la tension efficace, l'intensité efficace\ldots

\vspace{5mm}

\fcolorbox[gray]{0}{0.9}{\textbf{4.2. La valeur efficace d'une fonction développable en série de Fourier.}}\\[3mm]

Considérons une fonction de période $T$ dont le développement en série de Fourier est :
$$f(x)=a_0+\sum_{n=1}^{+\infty}\bigl(a_n\cos(n\omega x)+b_n\sin(n\omega x)\bigr)$$

Nous admettrons \textbf{la formule de Parseval} :

\begin{center}
\fbox{$\dfrac{1}{T}\displaystyle\int_\alpha^{\alpha+T}\bigl(f(x)\bigr)^2\,\mathrm{d}x
= a_0^2+\dfrac{1}{2}\displaystyle\sum_{n=1}^{+\infty}\bigl(a_n^2+b_n^2\bigr)$}
\end{center}

\vspace{5mm}

Si on note $f_{\text{eff}}$ la valeur efficace de la fonction $f$, ce résultat peut aussi s'écrire :

\begin{center}
\fbox{$f_{\text{eff}}^2=a_0^2+\dfrac{1}{2}\displaystyle\sum_{n=1}^{+\infty}\bigl(a_n^2+b_n^2\bigr)$}
\end{center}

\newpage
%%%%%%%%%%%%%%%%%%%%%%%%%%%%%%%%%%%%%%%%%%%%%%%%%%%%%%%%%%%%%
\fcolorbox[gray]{0}{0.9}{\textbf{Activité A5 -- Valeur efficace}}\\[2mm]

On considère la fonction $f$ définie par :
$\begin{cases}
    f(x) & \textrm{est paire},\quad f(x)\textrm{ est de période }2\pi \\
    f(x)=1 & \textrm{pour~~} x \in \left[0;\dfrac{\pi}{2}\right],\quad
    f(x)=0 \textrm{ pour~~} x \in \left]\dfrac{\pi}{2};\pi\right]
\end{cases}$

\begin{center}
\begin{tikzpicture}[x=1.0cm, y=1.8cm]
\def\xmin{-4.4} \def\xmax{6.4} \def\ymin{-1.4} \def\ymax{1.4}
\draw[xstep=0.25,ystep=0.1,very thin,orange!25] (\xmin,\ymin) grid (\xmax,\ymax);
\draw[xstep=1,ystep=0.5,thin,orange!55] (\xmin,\ymin) grid (\xmax,\ymax);
\draw[->,>=stealth',thick,blue!70!black] (\xmin,0)--(\xmax,0) node[right]{$x$};
\draw[->,>=stealth',thick,blue!70!black] (0,\ymin)--(0,\ymax) node[above]{$y$};
\foreach \n/\lab in {-4/{-2\pi},-3/{-\frac{3\pi}{2}},-2/{-\pi},-1/{-\frac{\pi}{2}},1/{\frac{\pi}{2}},2/{\pi},3/{\frac{3\pi}{2}},4/{2\pi},5/{\frac{5\pi}{2}},6/{3\pi}}
    \draw[blue!70!black] (\n,2pt)--(\n,-6pt) node[below]{\small$\lab$};
\foreach \y in {-1,-0.5,0.5,1}
    \draw[blue!70!black] (2pt,\y)--(-5pt,\y) node[left]{\small$\y$};
\begin{scope}
\clip (\xmin,\ymin) rectangle (\xmax,\ymax);
\foreach \c in {-4, 0, 4} {
    \draw[very thick, red!70!black] ({\c-1},1)--({\c+1},1);
    \draw[dashed, red!70!black] ({\c-1},0)--({\c-1},1);
    \draw[dashed, red!70!black] ({\c+1},0)--({\c+1},1);
}
\foreach \c in {-4, 0, 4} {
    \draw[very thick, red!70!black] ({\c+1},0)--({\c+3},0);
}
\end{scope}
\draw[black!60,thick] (\xmin,\ymin) rectangle (\xmax,\ymax);
\end{tikzpicture}
\end{center}

Nous avons obtenu le début du développement en série de Fourier :
$$g(x)=\frac{1}{2}+\frac{2}{\pi}\cos(x)-\frac{2}{3\pi}\cos(3x)$$
Calculer les valeurs efficaces $f_{\text{eff}}$ et $g_{\text{eff}}$ et donner la valeur approchée de $\dfrac{g_{\text{eff}}}{f_{\text{eff}}}$.

\textbf{Calcul de $f_{\text{eff}}$}

\cor{
$$f_{\text{eff}}^2 = \frac{1}{T}\int_0^{T}\bigl(f(x)\bigr)^2\,\mathrm{d}x
= \frac{1}{\pi}\int_0^{\pi}\bigl(f(x)\bigr)^2\,\mathrm{d}x \quad (f\text{ paire})
= \frac{1}{\pi}\int_0^{\pi/2}1\,\mathrm{d}x = \frac{1}{2}
\quad\Rightarrow\quad \boxed{f_{\text{eff}} = \frac{\sqrt{2}}{2} \approx 0{,}707}$$
}

\textbf{Calcul de $g_{\text{eff}}$ par la formule de Parseval} \quad ($a_0=\tfrac{1}{2}$, $a_1=\tfrac{2}{\pi}$, $a_3=\tfrac{-2}{3\pi}$, autres nuls)

\cor{
\begin{align*}
g_{\text{eff}}^2 &= a_0^2 + \frac{1}{2}(a_1^2+a_3^2)
= \frac{1}{4} + \frac{1}{2}\left(\frac{4}{\pi^2}+\frac{4}{9\pi^2}\right)
= \frac{1}{4} + \frac{2}{\pi^2}\cdot\frac{10}{9}
\approx 0{,}250 + 0{,}225 = 0{,}475 \\
&\Rightarrow\quad \boxed{g_{\text{eff}} \approx \sqrt{0{,}475} \approx 0{,}689}
\end{align*}
}

\textbf{Rapport $\dfrac{g_{\text{eff}}}{f_{\text{eff}}}$}

\cor{
$$\frac{g_{\text{eff}}}{f_{\text{eff}}} \approx \frac{0{,}689}{0{,}707} \approx \boxed{0{,}97}$$
La valeur efficace de $g$ (avec seulement 3 termes) représente environ $97\,\%$ de celle de $f$ : l'approximation est très bonne.
}

\vspace{3mm}
\textbf{4.3 Applications}\\
En physique cette formule permet de voir si l'utilisation des premiers harmoniques permet d'utiliser le début du développement en série de Fourier pour remplacer un signal \textbf{(voir A4)}.

\newpage
%%%%%%%%%%%%%%%%%%%%%%%%%%%%%%%%%%%%%%%%%%%%%%%%%%%%%%%%%%%%%%%%
\large
\begin{center}
\textbf{FICHE ESSENTIEL}\\[0.5cm]
\normalsize
\end{center}
\hrule
\vspace{8mm}
\begin{center}
\fbox{
\begin{minipage}{16.5cm}{
\textbf{Les séries numériques}\\

Les séries géométriques :\\
La série $\displaystyle \sum\limits_{n=0}^{+\infty}q^n$ converge si et seulement si $q \in ]-1;1[$.\\

Les séries de Riemann :\\
La série $\displaystyle \sum\limits_{n=1}^{+\infty}\dfrac{1}{n^\alpha}$ converge si et seulement si $\alpha>1$.\\[2mm]

\vspace{8mm}
\textbf{Le développement en série de Fourier d'une fonction}\\
Si $f$ est une fonction périodique de période $T$, on note \textbf{sa pulsation} $\omega = \dfrac{2\pi}{T}$\\
Son développement en série de Fourier est $S(x)=a_0+\displaystyle \sum\limits_{n=1}^{+\infty}(a_n\cos(n\omega x) + b_n\sin(n\omega x))$ avec :
\begin{center}
$a_0=\dfrac{1}{T}\displaystyle \int\limits_{\alpha}^{\alpha+T}f(x)\,\mathrm{d}x$ \quad (valeur moyenne de $f$)\\[4pt]
$a_n=\dfrac{2}{T}\displaystyle \int\limits_{\alpha}^{\alpha+T}f(x)\cos(n\omega x)\,\mathrm{d}x$ \qquad et \qquad $b_n=\dfrac{2}{T}\displaystyle \int\limits_{\alpha}^{\alpha+T}f(x)\sin(n\omega x)\,\mathrm{d}x$
\end{center}

\vspace{8mm}
\textbf{Cas particuliers}\\

Si $f$ est une fonction \textbf{impaire} : tous les coefficients $a_n$ (y compris $a_0$) sont nuls et
\begin{center}
\fbox{$b_n=\dfrac{4}{T}\displaystyle \int\limits_{0}^{T/2}f(x)\sin(n\omega x)\,\mathrm{d}x$}
\end{center}

Si $f$ est une fonction \textbf{paire} : tous les coefficients $b_n$ sont nuls et
\begin{center}
\fbox{$a_0=\dfrac{2}{T}\displaystyle \int\limits_{0}^{T/2}f(x)\,\mathrm{d}x$ \qquad et \qquad $a_n=\dfrac{4}{T}\displaystyle \int\limits_{0}^{T/2}f(x)\cos(n\omega x)\,\mathrm{d}x$}
\end{center}

\vspace{8mm}
\textbf{La formule de Parseval}\\
Si la fonction $f$ de période $T$ est développable en série de Fourier, sa valeur efficace est donnée par :
\begin{center}
\fbox{$f_{\text{eff}}^2=\dfrac{1}{T}\displaystyle \int\limits_{\alpha}^{\alpha+T}\bigl(f(x)\bigr)^2\,\mathrm{d}x=a_0^2+\dfrac{1}{2}\displaystyle \sum\limits_{n=1}^{+\infty}(a_n^2 + b_n^2)$}
\end{center}

\vspace{3mm}
}
\end{minipage}}
\end{center}

\newpage
%%%%%%%%%%%%%%%%%%%%%%%%%%%%%%%%%%%%%%%%%%%%%%%%%%%%%%%%%%%%%


%%%%%%%%%%%%%%%%%%%%%%%%%%%%%%%%%%%%%%%%%%%%%%%%%%%%%%%%%%%
\fcolorbox[gray]{0}{0.9}{\textbf{Extrait - Brevet de technicien supérieur 10 mai 2021 - groupement A}}\\[2mm]
\textbf{Exercice 1 \hfill 11 points}

\medskip

La communication par courants porteurs en ligne (ou CPL) permet de transmettre des informations en utilisant des conducteurs électriques en fonctionnement.

Le principe des CPL consiste à superposer au réseau électrique un signal de haute fréquence et de basse énergie. Ce deuxième signal se propage sur l'installation électrique et peut être reçu et décodé à distance.

\textbf{Partie A : Étude d'un signal}

\medskip

On s'intéresse à une tension périodique $U$. On note $T$ sa période.

\medskip

\begin{enumerate}
\item On donne ci-dessous une représentation graphique de $U$, exprimée en volt (V), en fonction de $t$ \textbf{exprimé en microseconde} ($\mu s$). Rappel : $1\,\mu s = 10^{-6}~s$.

\begin{center}
\begin{tikzpicture}[x=0.45cm, y=0.35cm]
\def\xmin{-11} \def\xmax{20} \def\ymin{-1.5} \def\ymax{14}
\draw[xstep=1,ystep=1,very thin,orange!25] (\xmin,\ymin) grid (\xmax,\ymax);
\draw[xstep=2,ystep=4,thin,orange!55] (\xmin,\ymin) grid (\xmax,\ymax);
\draw[->,>=stealth',thick,blue!70!black] (\xmin,0)--(\xmax,0) node[right]{\small$t~(\mu s)$};
\draw[->,>=stealth',thick,blue!70!black] (0,\ymin)--(0,\ymax) node[above]{\small$U~(\text{V})$};
\foreach \x in {-10,-8,-6,-4,-2,2,4,6,8,10,12,14,16,18}
    \draw[blue!70!black] (\x,2pt)--(\x,-5pt) node[below]{\tiny$\x$};
\foreach \y in {4,8,12}
    \draw[blue!70!black] (2pt,\y)--(-5pt,\y) node[left]{\small$\y$};
\draw[very thick,red!70!black] (-10,12)--(-2,12)--(-2,0)--(0,0);
\draw[very thick,red!70!black] (0,12)--(8,12)--(8,0)--(10,0);
\draw[very thick,red!70!black] (10,12)--(18,12)--(18,0)--(19,0);
\draw[very thick,red!70!black] (-10,0)--(-10,12);
\draw[very thick,red!70!black] (0,0)--(0,12);
\draw[very thick,red!70!black] (10,0)--(10,12);
\draw[black!60,thick] (\xmin,\ymin) rectangle (\xmax,\ymax);
\end{tikzpicture}
\end{center}

    \begin{enumerate}
        \item D'après la représentation graphique ci-dessus, quelle est la valeur de $T$ en microseconde ?\\
        \cor{D'après le graphique, on lit $\mathbf{T=10\,\mu s}$.}

        \item La fréquence d'un signal périodique, en Hz, suit la formule $f = \dfrac{1}{T}$ où $T$ est exprimé en seconde. Quelle est la fréquence de la tension périodique $U$ ?\\
        \cor{$$f = \frac{1}{T} = \frac{1}{10\times 10^{-6}} = \boxed{10^5\,\text{Hz} = 100\,\text{kHz}}$$}
    \end{enumerate}

\item On modélise l'évolution de $U$ (en volt) en fonction de $t$ (en microseconde) à l'aide d'une fonction numérique $f$ définie sur $\mathbb{R}$. Ainsi : $U = f(t)$.

On admet que la fonction $f$ est paire, périodique de période $T$, développable en série de Fourier et vérifie, pour tout réel $t$ :
\[f(t) = a_0 + \displaystyle\sum_{n=1}^{+\infty}\left(a_n\cos(n\omega t) + b_n\sin(n\omega t)\right)\quad\text{avec}\quad \omega = \dfrac{2\pi}{T}.\]

\newpage

    \begin{enumerate}
        \item Quelle est la valeur de $b_n$ pour tout entier $n\geqslant 1$ ? Justifier.\\
        \cor{$f$ est paire, donc tous les coefficients $b_n$ sont nuls : $\boxed{b_n = 0}$ pour tout $n\geqslant 1$.}

        \item $a_0 = \dfrac{1}{T}\displaystyle\int_0^{T} f(t)\:\mathrm{d}t$. Justifier que $a_0 = 9{,}6$.\\
        \cor{$$a_0 = \frac{1}{10}\int_0^{8}12\,\mathrm{d}t + \frac{1}{10}\int_8^{10}0\,\mathrm{d}t = \frac{12\times 8}{10} = \boxed{9{,}6\,\text{V}}$$}
    \end{enumerate}

\item La valeur efficace de $U$, notée $U_{\text{eff}}$, vérifie : $U^2_{\text{eff}} = \dfrac{1}{T}\displaystyle\int_0^{T}[f(t)]^2\:\mathrm{d}t$. Montrer que : $U_{\text{eff}} = 10{,}7$~V.\\
\cor{$$U_{\text{eff}}^2 = \frac{1}{10}\int_0^{8}144\,\mathrm{d}t = \frac{144\times 8}{10} = 115{,}2 \qquad\Rightarrow\qquad U_{\text{eff}} = \sqrt{115{,}2} \approx \boxed{10{,}7\,\text{V}}$$}

\item Le signal correspondant à la tension $U$ est envoyé sur une ligne moyenne tension transportant une tension efficace de $M = 20\,000$~V.

Le taux de distorsion harmonique par rapport au fondamental, noté $T_F$, est donné par la formule :
\[T_F = \sqrt{\dfrac{U^2_{\text{eff}} - a_0^2}{M^2}}.\]
On considère qu'un CPL n'a pas d'incidence sur le réseau si $T_F$ est inférieur à 0,1\,\%. Le CPL étudié dans la partie A a-t-il une incidence sur le réseau ?\\
\cor{
$$T_F = \sqrt{\frac{115{,}2 - (9{,}6)^2}{(20\,000)^2}} = \sqrt{\frac{23{,}04}{4\times10^8}} = \sqrt{5{,}76\times10^{-8}} = 2{,}4\times10^{-4} = 0{,}024\,\%$$
Comme $T_F = 0{,}024\,\% < 0{,}1\,\%$, le CPL \textbf{n'a pas d'incidence} sur le réseau.
}
\end{enumerate}

\newpage
%%%%%%%%%%%%%%%%%%%%%%%%%%%%%%%%%%%%%%%%%%%%%%%%%%%%%%%%%%%%%
\fcolorbox[gray]{0}{0.9}{\textbf{Extrait - Brevet de technicien supérieur septembre 2020 - groupement A}}\\[2mm]

Le moteur triphasé d'une voiture électrique est alimenté par un onduleur pour modifier sa vitesse.

La tension simple à la sortie de cet onduleur dépend de $\omega t$ où $\omega$ est la pulsation (en rad/s) et $t$ le temps (en s). On l'exprime sous la forme $v(\omega t)$, où $v$ est une fonction paire, périodique de période $T = 2\pi$, et définie pour tout réel $x$ appartenant à l'intervalle $[0;\pi]$ par :
\[\left\{\begin{array}{lcll}
v(x) &=& 300  &\text{si}\quad 0 \leqslant x < \dfrac{\pi}{6} \\[6pt]
v(x) &=& 150  &\text{si}\quad \dfrac{\pi}{6} \leqslant x \leqslant \dfrac{\pi}{2}\\[6pt]
v(x) &=& -150 &\text{si}\quad \dfrac{\pi}{2} < x < \dfrac{5\pi}{6}\\[6pt]
v(x) &=& -300 &\text{si}\quad \dfrac{5\pi}{6} \leqslant x \leqslant \pi
\end{array}\right.\]

\medskip

\begin{enumerate}
\item Compléter la représentation graphique de la fonction $v$ sur l'intervalle $\left[-\dfrac{3\pi}{2};\dfrac{3\pi}{2}\right]$.
\begin{center}
\begin{tikzpicture}[x=0.7cm, y=0.38cm]
\def\xmin{-9.5} \def\xmax{9.5} \def\ymin{-7.5} \def\ymax{7.5}
\draw[xstep=1,ystep=1,very thin,orange!25] (\xmin,\ymin) grid (\xmax,\ymax);
\draw[xstep=3,ystep=2,thin,orange!55] (\xmin,\ymin) grid (\xmax,\ymax);
\draw[->,>=stealth',thick,blue!70!black] (\xmin,0)--(\xmax,0) node[right]{\small$x$};
\draw[->,>=stealth',thick,blue!70!black] (0,\ymin)--(0,\ymax) node[above]{\small$v$};
\foreach \n/\lab in {-9/{-\frac{3\pi}{2}},-6/{-\pi},-3/{-\frac{\pi}{2}},3/{\frac{\pi}{2}},6/{\pi},9/{\frac{3\pi}{2}}}
    \draw[blue!70!black] (\n,3pt)--(\n,-5pt) node[below]{\small$\lab$};
\foreach \y/\lab in {-6/{-300},-3/{-150},3/{150},6/{300}}
    \draw[blue!70!black] (2pt,\y)--(-5pt,\y) node[left]{\small$\lab$};
% Courbe complète (prof) sur [-9.5, 9.5]
% Parité + périodicité : paires de valeurs en TikZ (1 unité = π/6, 50V=1 unité)
\draw[very thick,red!70!black]
    (-9.5,3)--(-9,3)--(-9,-3)--(-7,-3)--(-7,-6)--(-5,-6)--(-5,-3)--(-3,-3)--(-3,3)--(-1,3)--(-1,6)--
    (0,6)--(1,6)--(1,3)--(3,3)--(3,-3)--(5,-3)--(5,-6)--
    (6,-6)--(7,-6)--(7,-3)--(9,-3)--(9,3)--(9.5,3);
\draw[black!60,thick] (\xmin,\ymin) rectangle (\xmax,\ymax);
\end{tikzpicture}
\end{center}
\cor{La fonction $v$ est paire (symétrie par rapport à l'axe des ordonnées) et $2\pi$-périodique : on complète par réflexion puis par translation.}

On admet que la fonction $v$ est développable en série de Fourier et que, pour tout réel $x$, on a :
\[v(x) = a_0 + \displaystyle\sum_{n=1}^{+\infty}\left(a_n\cos(nx) + b_n\sin(nx)\right).\]

\item Justifier que $b_n = 0$ pour tout entier $n\geqslant 1$.\\
\cor{$v$ est paire, donc $b_n=0$ pour tout $n\geqslant 1$.}

\item On rappelle que : $a_0 = \dfrac{2}{T}\displaystyle\int_0^{T/2} v(x)\:\mathrm{d}x$. Montrer que : $a_0 = 0$.\\
\cor{
$$a_0 = \frac{1}{\pi}\int_0^{\pi}v(x)\,\mathrm{d}x = \frac{1}{\pi}\!\left[300\cdot\frac{\pi}{6}+150\left(\frac{\pi}{2}-\frac{\pi}{6}\right)-150\left(\frac{5\pi}{6}-\frac{\pi}{2}\right)-300\left(\pi-\frac{5\pi}{6}\right)\right]$$
$$= \frac{1}{\pi}\bigl[50\pi+50\pi-50\pi-50\pi\bigr] = \boxed{0}$$
}

\item On donne, pour tout entier $n\geqslant 1$ : $a_n = \dfrac{4}{T}\displaystyle\int_0^{T/2}v(x)\cos(nx)\:\mathrm{d}x$.
    \begin{enumerate}
        \item On pose : $I_n = \displaystyle\int_0^{\pi/6}300\cos(nx)\:\mathrm{d}x$ et $J_n = \displaystyle\int_{5\pi/6}^{\pi}\left[-300\cos(nx)\right]\:\mathrm{d}x$. Calculer $I_n$ et $J_n$ en fonction de $n$.\\
        \cor{
        $$I_n = 300\left[\frac{\sin(nx)}{n}\right]_0^{\pi/6} = \frac{300}{n}\sin\!\left(\frac{n\pi}{6}\right) \qquad
        J_n = \frac{300}{n}\sin\!\left(\frac{5n\pi}{6}\right)$$
        }

        \item Écrire les intégrales qu'il reste à calculer pour obtenir l'expression de $a_n$ en fonction de $n$. \textbf{On ne demande pas de les calculer.}\\
        \cor{
        $$K_n = \int_{\pi/6}^{\pi/2}150\cos(nx)\,\mathrm{d}x \qquad\text{et}\qquad L_n = \int_{\pi/2}^{5\pi/6}(-150)\cos(nx)\,\mathrm{d}x$$
        Puis : $a_n = \dfrac{2}{\pi}(I_n+K_n+L_n+J_n)$.
        }
    \end{enumerate}

\item On admet que, pour tout entier $n\geqslant 1$ : $a_n = \dfrac{300}{n\pi}\!\left[\sin\!\left(\dfrac{n\pi}{6}\right)+2\sin\!\left(\dfrac{n\pi}{2}\right)+\sin\!\left(\dfrac{5n\pi}{6}\right)\right]$.

Compléter le tableau ci-dessous :
\begin{center}
\begin{tabularx}{\linewidth}{|m{4.2cm}|*{8}{>{\centering\arraybackslash}X|}}
\hline
Valeurs de $n$ & 0 & 1 & 2 & 3 & 4 & 5 & 6 & 7\\ \hline
Valeurs approchées de $a_n$ arrondies au dixième. & 0 & 286,5 & 0 & 0 & \cor{0} & \cor{57,3} & \cor{0} & \cor{$-40{,}9$}\\ \hline
\end{tabularx}
\end{center}
\cor{\textit{$n=5$} : $\frac{300}{5\pi}[\frac{1}{2}+2+\frac{1}{2}]=\frac{180}{\pi}\approx 57{,}3$ \qquad \textit{$n=7$} : $\frac{300}{7\pi}[-\frac{1}{2}-2-\frac{1}{2}]=\frac{-900}{7\pi}\approx -40{,}9$}

\item On évalue la pollution harmonique de deux façons : de manière qualitative, en observant la présence d'harmoniques sur le spectre des amplitudes, et de manière quantitative en calculant le taux de distorsion harmonique par rapport au fondamental.
    \begin{enumerate}
        \item Compléter sur le graphique ci-dessous le spectre des amplitudes $A_n$ associé à la fonction $v$ pour $n$ allant de $0$ à $7$.\\
        On rappelle que pour tout entier $n\geqslant 1$ : $A_n = \sqrt{a_n^2+b_n^2}$.
        \begin{center}
        \begin{tikzpicture}[x=1.6cm, y=0.013cm]
\def\xmin{-0.5} \def\xmax{7.8} \def\ymin{-20} \def\ymax{320}
\draw[xstep=1,ystep=50,very thin,orange!25] (\xmin,0) grid (\xmax,\ymax);
\draw[xstep=1,ystep=100,thin,orange!55] (\xmin,0) grid (\xmax,\ymax);
\draw[->,>=stealth',thick,blue!70!black] (\xmin,0)--(\xmax,0) node[right]{\small$n$};
\draw[->,>=stealth',thick,blue!70!black] (0,\ymin)--(0,\ymax) node[above]{\small$A_n$~(V)};
\foreach \n in {0,1,2,3,4,5,6,7}
    \draw[blue!70!black] (\n,3pt)--(\n,-8pt) node[below]{\small$\n$};
\foreach \y in {50,100,150,200,250,300}
    \draw[blue!70!black] (2pt,\y)--(-6pt,\y) node[left]{\small$\y$};
% A_1 ≈ 286.5 (donné)
\fill[red!60!white] (0.7,0) rectangle (1.3,286.5);
\draw[red!70!black,thick] (0.7,0) rectangle (1.3,286.5);
\node[above,red!70!black,font=\small] at (1,286.5) {$\approx287$};
% A_5 ≈ 57.3 (complété)
\fill[red!60!white] (4.7,0) rectangle (5.3,57.3);
\draw[red!70!black,thick] (4.7,0) rectangle (5.3,57.3);
\node[above,red!70!black,font=\small] at (5,57.3) {$\approx57$};
% A_7 ≈ 40.9 (complété)
\fill[red!60!white] (6.7,0) rectangle (7.3,40.9);
\draw[red!70!black,thick] (6.7,0) rectangle (7.3,40.9);
\node[above,red!70!black,font=\small] at (7,40.9) {$\approx41$};
% A=0 pour n=0,2,3,4,6
\foreach \n in {0,2,3,4,6}
    \draw[blue!50!black,thick] (\n-0.3,0)--(\n+0.3,0);
\draw[black!60,thick] (\xmin,\ymin) rectangle (\xmax,\ymax);
        \end{tikzpicture}
        \end{center}
        \cor{$A_0=0$,\; $A_1\approx286{,}5$,\; $A_2=A_3=A_4=A_6=0$,\; $A_5\approx57{,}3$,\; $A_7\approx40{,}9$.}

        \item On s'intéresse au taux de distorsion harmonique par rapport au fondamental.

On obtient une valeur approchée significative de ce taux en calculant le nombre T donné par :
\[\text{T} = \dfrac{\sqrt{a_2^2+a_3^2+a_4^2+a_5^2+a_6^2+a_7^2}}{a_1}.\]
Sachant que ce taux doit être inférieur à 10\,\%, la commande de l'onduleur du moteur est-elle adaptée ?\\
        \cor{
        $$\text{T} = \frac{\sqrt{\left(\frac{180}{\pi}\right)^2+\left(\frac{900}{7\pi}\right)^2}}{\frac{900}{\pi}}
        = \frac{1}{900}\sqrt{180^2+\left(\frac{900}{7}\right)^2}
        \approx \frac{221{,}2}{900} \approx 0{,}246$$
        $\text{T}\approx24{,}6\,\%>10\,\%$ : la commande de l'onduleur \textbf{n'est pas adaptée} (taux trop élevé).
        }
    \end{enumerate}
\end{enumerate}

\end{document}
